
\section{Ontología dual: materia positiva y arquitectura negativa}

Sea \(D\subset\R^3\) el dominio de construcción y \(M(t)\subseteq D\) el dominio material. El espacio libre es

\begin{equation}
V(t)=D\setminus M(t).
\end{equation}

IRL modela simultáneamente

\begin{equation}
\boxed{
\mathcal A(t)
=
\left(
\Kp(t),
\Kn(t)
\right),
}
\end{equation}

donde \(\Kp\) es el complejo material y \(\Kn\) el complejo arquitectónico negativo.

\begin{definition}[Objeto negativo arquitectónico]
Un objeto negativo \(N_j\) se representa mediante
\begin{equation}
N_j
=
\left(
V_j,
\partial_sV_j,
P_j,
\mathcal H_j,
m_j^\ominus,
\eta_j,
I_j,
\mathcal E_j
\right),
\end{equation}
donde \(V_j\) es volumen, \(\partial_sV_j\) frontera estructural, \(P_j\) portales, \(\mathcal H_j\) firma topológica, \(m_j^\ominus\) déficit de masa de referencia, \(\eta_j\) grado de realización de frontera, \(I_j\) estado de identidad y \(\mathcal E_j\) evidencia.
\end{definition}

Una cámara no es sólo aire: es una región cuya no-ocupación tuvo que preservarse mientras se construían sus límites y dependencias.

\section{Masa ausente contrafactual}

Sea \(\rho_{\mathrm{ref}}(\mathbf x)\) una densidad de referencia explícita.

\begin{definition}[Masa ausente contrafactual]
\begin{equation}
\boxed{
m_C^\ominus
=
-\int_{V_C}
\rho_{\mathrm{ref}}(\mathbf x)\,dV.
}
\end{equation}
\end{definition}

\(m_C^\ominus\) no es masa física negativa ni produce antigravedad. Es un déficit firmado respecto de una hipótesis de relleno.

Definimos también

\begin{equation}
\delta\rho(\mathbf x)
=
\rho_{\mathrm{obs}}(\mathbf x)
-
\rho_{\mathrm{ref}}(\mathbf x).
\end{equation}

La muografía restringe integrales de densidad a lo largo de trayectorias,

\begin{equation}
\Sigma_\mu(\gamma)
=
\int_\gamma \rho(\mathbf x)\,ds,
\end{equation}

por lo que una anomalía de densidad puede incorporarse al modelo sin promoverla automáticamente a cámara arquitectónica \cite{morishima2017,procureur2023}.

\section{Genealogía del vacío y campos de nacimiento}

Una cámara puede atravesar al menos cinco eventos:

\begin{enumerate}[label=(\roman*)]
\item reserva del volumen futuro \(t_R\);
\item emergencia de fronteras \(t_B\);
\item cierre espacial \(t_C\);
\item completitud estructural \(t_S\);
\item encapsulamiento \(t_E\).
\end{enumerate}

Con

\begin{equation}
t_R\le t_B\le t_C\le t_S\le t_E.
\end{equation}

Se proponen dos campos:

\begin{equation}
\boxed{
T_b^+(\mathbf x)
=
\text{instante de incorporación material}
}
\end{equation}

y

\begin{equation}
\boxed{
T_b^-(\mathbf x)
=
\text{instante de reserva del espacio negativo}.
}
\end{equation}

La consecuencia es

\begin{equation}
\boxed{
\text{vacío geométrico}
\neq
\text{identidad topológica}
\neq
\text{identidad arquitectónica}.
}
\end{equation}

\section{Bifiltración del espacio negativo}

La conectividad ordinaria no distingue una cámara de un corredor conectado.

Sea

\begin{equation}
d_M(\mathbf x,\tau)
=
\operatorname{dist}
(\mathbf x,M_\tau).
\end{equation}

Para un radio de despeje \(r\),

\begin{equation}
\boxed{
V_{\tau,r}
=
\left\{
\mathbf x\in D\setminus M_\tau:
d_M(\mathbf x,\tau)\ge r
\right\}.
}
\end{equation}

El mapa

\begin{equation}
(\tau,r)\mapsto V_{\tau,r}
\end{equation}

forma una bifiltración propuesta. A \(\tau\) fijo, aumentar \(r\) elimina conexiones estrechas antes que grandes recintos. A \(r\) fijo, aumentar \(\tau\) recorre estados inversos.

Pueden analizarse

\begin{equation}
H_k(V_{\tau,r}),
\qquad
\beta_k(\tau,r),
\end{equation}

además de merge trees, grafos de Reeb y persistencia multiparamétrica \cite{edelsbrunner2010,carlsson2009,carlssonzomorodian2009}.

\begin{conjecture}[Cronología topológica]
Eventos persistentes de espacio negativo pueden imponer cotas cronológicas no recuperables únicamente desde elevación.
\end{conjecture}
